\section{The eight small quadratic twists of the mixed cubic family}
\label{sec:mixed-cubic-quadratic-window}

\paragraph{The preceding original-coefficient remainder.}
Retain every datum, page, prescribed pair, host, original coefficient,
lattice, normalization and full local condition of AC-G206-window.
In particular $h\geq2$ and $T\geq2$ are integers, $N=T-1$, and
\[
 \begin{gathered}
 H_u=\mathbb Q(\zeta_{2^{h+2}},\zeta_7,\sqrt{-q},\sqrt{-d},
                  \omega,\sqrt[3]{269}),\qquad K=\mathbb Q(\omega),\\
 u^{2+8\varepsilon}\leq2^h\leq2u^{2+8\varepsilon},\qquad
 T=\lfloor u^{399}\rfloor,\qquad0<\varepsilon<1/256.
 \end{gathered}\tag{MQW.1}
\]
Let $\chi$ be either faithful character of
$\operatorname{Gal}(K(\sqrt[3]{269})/K)$.
Keep exactly $Z_u,J_{u,0},J_{u,\rm mix},V_u,Y_u$,
$C_5^{\rm all}(u),U_5^{\rm all}(u),C_7^{\rm all}(u),U_7^{\rm all}(u)$
and $M(u)$ from (C5A.2)--(C5A.6).
The coefficient-prime set $\mathscr P_{\rm cub}$ and
$J_{\rm cub}(T)$ are precisely (MCW.2), equivalently (MCQ.3):
\[
 \begin{split}
 \mathscr P_{\rm cub}=\{\ell\text{ prime}:{}&\ell>3,\quad
    \ell\ne7,269,\quad\ell\equiv1\pmod3,\\
    &269^{(\ell-1)/3}\not\equiv1\pmod\ell\},\\
 Q(u)=\{(\ell,\chi^j\delta,e):{}&\ell\in\mathscr P_{\rm cub},\quad
       j=1,2,\quad\delta\in\{\chi_{-4},\chi_8,\chi_{-8}\},\\
       &1\leq e<T,\quad(\ell-1)\mid e\},\\
 M_8(u)&=M(u)\cup Q(u),\\
 J_{\rm cub}(T)&=\sum_{\substack{\ell\in\mathscr P_{\rm cub}\\\ell\leq T}}
                         \left\lfloor\frac{N}{\ell-1}\right\rfloor.
 \end{split}\tag{MQW.2}
\]
Every undisplayed conductor-seven and auxiliary $q,d$ factor
is trivial. The three displayed quadratic factors are the original
two-power characters; no finite character at $\ell$ is added.

Define the explicit remainder
\[
 \begin{split}
 \mathcal R_{399}^{\rm match,cub8,5all}(u)={}&
  \sum_{(\ell,\eta,e)\in J_{u,\rm mix}\setminus M_8(u)}
                    \frac{\log\#\mathcal S_{\ell,\eta}(e)}{b_\ell}\\
 &+\sum_{(\ell,\eta,e)\in
       J_{u,0}\setminus(V_u\cup Y_u\cup C_7^{\rm all}(u)
                                       \cup C_5^{\rm all}(u))}
       \frac{\log\#\mathcal S^{\mathbb Q}_{\ell,
                            \psi_{\eta,e}^{\rm nc}}(e)}{b_\ell}.
 \end{split}\tag{MQW.3}
\]
These are the same finite nonnegative strict-group sums as in
(C5A.3). The rational noncritical sum is unchanged.

\begin{theorem}[AC-G208-window: quadratic twists on the same mixed window]
\label{AC-G208-window}
Every triple of $Q(u)$ has zero full original $K$-strict group.
The complete old-overlap statement and pointwise identity are
\[
 \begin{gathered}
 Q(u)\subset J_{u,\rm mix}\setminus M(u),\\
 Q(u)\cap\bigl(Z_u\cup V_u\cup Y_u\cup C_7^{\rm all}(u)
                              \cup C_5^{\rm all}(u)\cup M(u)\bigr)
                                      =\varnothing,\\
 M_8(u)=M(u)\sqcup Q(u),\qquad
 \mathcal R_{399}^{\rm match,cub8,5all}(u)
                 =\mathcal R_{399}^{\rm match,cub,5all}(u).
 \end{gathered}\tag{MQW.4}
\]
The exact counts are
\[
 \begin{gathered}
 |Q(u)|=6J_{\rm cub}(T),\qquad |M_8(u)|=8J_{\rm cub}(T),\\
 |Q(u)|\geq6\left\lfloor N/12\right\rfloor
                       +6\left\lfloor N/18\right\rfloor.
 \end{gathered}\tag{MQW.5}
\]
Thus exactly $6J_{\rm cub}(T)$ additional mixed terms are absent
relative to AC-G206-window, with a proper indexing reduction
for $T\geq13$. The full disjoint represented-triple omission is
\[
 \begin{gathered}
 Z_u\sqcup V_u\sqcup Y_u\sqcup U_7^{\rm all}(u)
                   \sqcup U_5^{\rm all}(u)\sqcup M_8(u),\\
 |Z_u|+|V_u|+|Y_u|+47N-\lfloor N/4\rfloor+8J_{\rm cub}(T).
 \end{gathered}\tag{MQW.6}
\]
The second line is its cardinality. With exactly the preceding
constants and costs,
\[
 \begin{aligned}
 |I_{399}-\mathcal R_{399}^{\rm match,cub8,5all}|&\leq
          \mathcal E_u+\mathcal L_u+\mathcal A_u+C_2\mathcal P_u,\\
 |I_{399}^{\rm all7,2}-\mathcal R_{399}^{\rm match,cub8,5all}|&\leq
          \mathcal E_u+\mathcal L_u+C_2\mathcal P_u.
 \end{aligned}\tag{MQW.7}
\]
No extra local, torsion, critical or family-count cost occurs.
The envelope remains $O_\varepsilon(u^{806+16\varepsilon}(\log u)^2)$.
Equal explicit remainder premises use the identical window,
threshold and slack; the original full index premise follows
with the same threshold increase and half-slack argument.
\end{theorem}

\begin{proof}
\proofstep{1. Actual original characters and literal strict-zero input.}
The field in (MQW.1) contains $\mathbb Q(\zeta_8)$.
Its three quadratic characters $\chi_{-4},\chi_8,\chi_{-8}$
are the characters of the biquadratic field $\mathbb Q(i,\sqrt2)$.
They remain distinct and nontrivial on $G_K$: none of that
field's three quadratic subfields is $K$.
The product decomposition (FC.1) makes the Kummer $C_3$
independent of these characters and of the conductor-seven
and auxiliary factors. Consequently the six twisted labels
are distinct, and none equals an untwisted $\chi^j$.
This also prevents confusing the Kummer factor with a rational
conductor-seven character of order three.

Their orders divide six and hence $3\cdot2^h$.
Use the prescribed original completion $O_\ell$ of
$\mathbb Z[\zeta_{3\cdot2^h}]$, its fraction field $E_\ell$,
uniformizer and normalizing rank $b_\ell$.
The literal coefficient in (FC.2) is canonically
\[
 (O_\ell\otimes_{\mathbb Z_\ell}\mathbb Q_\ell/\mathbb Z_\ell)
                         (\chi^j\delta)(-e)
       =(E_\ell/O_\ell)(\chi^j\delta c_\ell^{-e}).
 \tag{MQW.8}
\]
Apply AC-G207 (MCQ.4)--(MCQ.5) with $\kappa=1$ and the
same $\ell,e,j,\delta$. Its coefficient-prime and weight
conditions are exactly (MQW.2). The actual-host inflation
in G207 Step6 and (FC.2) identifies the required full strict
kernel with a subgroup of the zero group proved there.
Both split coefficient-prime completions and every away-inertia
condition, including those at $2,3,269$, are retained.
The larger original coefficient ring and its embedding are
unchanged even though the finite character values already
lie in a smaller field.

The class input to this application is exactly (MCQ.2):
the three degree-six groups are quotients of $C_7,C_{63},C_{42}$.
Their prime divisors impose no further restriction on
$\mathscr P_{\rm cub}$. The fixed-transposition norm argument
in G207 Step3 kills only the required nontrivial cubic
components of the degree-twelve class groups.
It does not assert that those entire class groups vanish.
The two local residual Frobenius values remain nontrivial
after multiplication by a quadratic sign, and G207's
original-uniformizer lifting retains all full local conditions.
No new class computation, unit-basis assumption or local-rank
premise is introduced by this composition.

The finite modules in (MQW.8) have the retained finite-cover
interpretation on
$\operatorname{Spec}\mathcal O_K[1/(2\cdot3\cdot269\cdot\ell)]$
from G207's dictionary. This is the same full-completion and
inertia kernel, with the retained SGA1 Expose V, sections4 and7
interpretation; no new motivic realization is claimed.

\proofstep{2. Exhaustive old intersections.}
Every new label has nontrivial Kummer cubic factor.
The old-overlap audit (MCW.4) and its Step2 identify all
old mixed omissions in $Z_u$ at coefficient $269$.
That coefficient is excluded in (MQW.2); the fixed two-
and three-coefficient families are excluded as well.
The other old $Z_u$ families, including (OPW.3) and (I2W.3),
have trivial Kummer factor and cannot contain a new row.
The sets $V_u,Y_u,C_7^{\rm all}(u),C_5^{\rm all}(u)$ also
have trivial Kummer factor. This distinction persists when
$\ell$ equals an auxiliary prime $q$ or $d$: an old rational
Legendre row in $Y_u$ still has trivial Kummer factor.
Thus every new triple lies in the retained $J_{u,\rm mix}$.

The only subsequent mixed omission is $M(u)$, whose quadratic
two-power factor is trivial. Every label in $Q(u)$ has one
of the three nontrivial quadratic factors, so Step1 shows
$M(u)\cap Q(u)=\varnothing$ at every prime and weight.
The G205/G206 completion changed only the rational noncritical
sum. It introduced no additional mixed omission.
This proves every disjointness assertion in (MQW.4) and
the union in (MQW.6), retaining all other earlier families.

\proofstep{3. Exact counts and pointwise mixed deletion.}
For each eligible prime the positive weights below the integer
$T$ are $(\ell-1),2(\ell-1),\ldots$, so their number is
$\lfloor N/(\ell-1)\rfloor$. A nonzero contribution implies
$\ell\leq T$. This is a consequence of the weight condition,
not a new upper bound in the unrestricted prime domain (TC.1).
There are six distinct new labels at each such weight and
two old untwisted labels. This proves the first two counts
in (MQW.5). The previously checked primes13 and19 in (MCW.5)
give its lower bound. In particular $(13,\chi\chi_{-4},12)$
is a new triple whenever $T\geq13$.
Combining $|U_7^{\rm all}|=23N$, $|U_5^{\rm all}|=24N-\lfloor N/4\rfloor$
and $|M_8|=8J_{\rm cub}(T)$ proves the cardinality in (MQW.6).

Step1 makes each new full mixed summand equal to
$\log1/b_\ell=0$. Delete these summands from (C5A.3), or
apply the arbitrary-retained-set identity (MCQ.8) after all
preceding omissions. The result is precisely (MQW.3), proving
the pointwise equality in (MQW.4).
The whole rational noncritical sum and its critical comparison
remain unchanged. There is no new critical partner for a
deleted full mixed group. Equality of sums does not assert a
strictly smaller numerical value or a zero individual derived index.

\proofstep{4. One inherited cost on the original common window.}
Substitute (MQW.4) into (C5A.7) to obtain (MQW.7).
This retains one full (FL.4) local majorant, one (LT.2) torsion
aggregate, one existing critical aggregate and the old G189
cost only against the full index. No multiplier for six,
$J_{\rm cub}(T)$, the number of primes or weights, or the
three proof fields is introduced. The degrees
$D_u=144\cdot2^h$ and $d_u=96\cdot2^h$ keep their separate roles.

Fix the same original datum, page, $0<\varepsilon<1/256$,
parameters $J,X_0,s$, prescribed pairs and actual window
$X^{1/J}\leq W$, $W^2\leq X$ of (TC.3).
Equal explicit remainders use precisely that window with
the same threshold and slack. For the full index retain the
preceding uniform $U_0$, increase the threshold to
$\max(X_0,U_0^J)$, and use the same half-slack bound on this
identical $W$. No independent existential windows are intersected.
Every simultaneous original field, coefficient, lattice, finite
condition, actual source rank $R.t$, page and pair remains.
\end{proof}

\paragraph{Scope and chronology.}
The eight labels refer to the two Kummer characters times
the four quadratic factors $1,\chi_{-4},\chi_8,\chi_{-8}$,
with trivial conductor-seven and auxiliary factors, at precisely
the displayed primes and weights. Every other mixed term and
every previously retained noncritical term remains in (MQW.3).
The remaining numerical bound, mixed theta/Selmer realization
and H04.4 remain open. The G207 arbitrary-set identity and
this complete original-window application are assessed together
at their actual mathematical completion times; packaging supplies
no separate movement or later banking. No new question is created.
